%==========================================================
% Chapter: Linear Classification in PyTorch (Student-Friendly)
% NOTE: This is LaTeX-only, with colors and improved formatting.
%==========================================================

\chapter{Linear Classification in PyTorch: From Geometry to Code (with Logistic Regression)}

\begin{center}
{\Large \color{blue!70!black}\textbf{Goal:}} \; Build and understand a binary classifier in PyTorch, \\
from the decision boundary equation to training, evaluation, and saving the model.
\end{center}

\vspace{0.5em}

%---------------------------
% Recommended packages (place in preamble)
%---------------------------
% \usepackage{amsmath, amssymb, mathtools}
% \usepackage{xcolor}
% \usepackage{hyperref}
% \usepackage{tcolorbox}
% \usepackage{listings}
% \usepackage{graphicx}
% \usepackage{enumitem}
% \usepackage{microtype}

%---------------------------
% Listings settings (place in preamble)
%---------------------------
% \lstset{
%   basicstyle=\ttfamily\small,
%   keywordstyle=\color{blue!70!black}\bfseries,
%   commentstyle=\color{green!50!black},
%   stringstyle=\color{orange!70!black},
%   showstringspaces=false,
%   frame=single,
%   rulecolor=\color{black!20},
%   breaklines=true,
%   columns=fullflexible
% }

%---------------------------
% Color boxes (place in preamble)
%---------------------------
% \tcbset{
%   colframe=blue!50!black,
%   colback=blue!2,
%   arc=2mm,
%   boxrule=0.6pt
% }

\section{The Big Picture: What Your ML Script Is Really Doing}

Most machine learning scripts follow this pattern:

\begin{tcolorbox}[title=\textbf{\color{blue!70!black}Standard workflow (high-level)}]
\begin{enumerate}[leftmargin=*]
  \item Load data: features $X$ and labels $y$
  \item Create a model
  \item Train the model (fit)
  \item Make predictions
  \item Evaluate performance
\end{enumerate}
\end{tcolorbox}

In scikit-learn, this is hidden inside functions like \texttt{fit}, \texttt{predict}, and \texttt{score}.
In PyTorch, we write the training loop ourselves, which is the best way to actually understand learning.

\section{Binary Classification: Accuracy, Error, and What We Care About}

\subsection{Targets and predictions}
Binary classification means:
\[
y \in \{0,1\}.
\]
Examples:
\begin{itemize}
  \item spam (1) vs not spam (0)
  \item malignant (1) vs benign (0)
\end{itemize}

\subsection{Accuracy}
\[
\text{Accuracy} = \frac{\#\text{correct}}{\#\text{total}}
= \frac{1}{N}\sum_{i=1}^{N}\mathbb{1}[\hat{y}_i = y_i].
\]

\subsection{Classification error}
\[
\text{Error} = 1 - \text{Accuracy}.
\]

\begin{tcolorbox}[colframe=green!50!black,colback=green!3,title=\textbf{\color{green!40!black}Student check: train vs test accuracy}]
\textbf{Train accuracy} measures performance on data the model has already seen. \\
\textbf{Test accuracy} measures performance on data the model has \textit{not} seen. \\
\textbf{\color{red!70!black}{Test accuracy is what you should trust more.}}
\end{tcolorbox}

\section{Geometry First: The Decision Boundary}

\subsection{A line / hyperplane}
In 2D, a separating line can be written as:
\[
w_1x_1 + w_2x_2 + b = 0.
\]
In $D$ dimensions, we use vectors:
\[
x \in \mathbb{R}^D,\quad w \in \mathbb{R}^D,
\]
and the compact form:
\[
\underbrace{a}_{\text{\color{blue!70!black}{activation / logit}}} = w^\top x + b.
\]

\begin{tcolorbox}[colframe=blue!50!black,colback=blue!2,title=\textbf{\color{blue!70!black}Interpretation}]
\begin{itemize}[leftmargin=*]
  \item $w$ controls \textbf{how strongly} each feature affects the decision.
  \item $b$ shifts the boundary left/right (or more generally: shifts the hyperplane).
  \item The value $a = w^\top x + b$ tells you \textbf{which side} of the boundary $x$ lies on.
\end{itemize}
\end{tcolorbox}

\subsection{Prediction from geometry (sign rule)}
If $a$ is positive, you are on one side; if negative, you are on the other:
\[
\hat{y} =
\begin{cases}
1 & \text{if } a>0,\\
0 & \text{otherwise}.
\end{cases}
\]

\section{From Hard Decisions to Probabilities: The Sigmoid}

The hard step rule is not smooth, and deep learning likes smooth functions.
So we map $a$ to a probability using the \textbf{sigmoid}:

\[
\sigma(a)=\frac{1}{1+e^{-a}},
\qquad
p(y=1\mid x)=\sigma(w^\top x+b).
\]

\subsection{Thresholding probabilities}
A common rule:
\[
\hat{y} =
\begin{cases}
1 & \text{if } p\ge 0.5,\\
0 & \text{otherwise}.
\end{cases}
\]

\begin{tcolorbox}[colframe=purple!60!black,colback=purple!3,title=\textbf{\color{purple!70!black}Key equivalence (very important!)}]
\[
\sigma(a)\ge 0.5 \;\Longleftrightarrow\; a\ge 0.
\]
So you can classify using either probabilities \textbf{or} logits (activations).
\end{tcolorbox}

\section{Training: Why Loss Functions Change for Classification}

\subsection{Why not mean squared error?}
Regression predicts real numbers, so MSE makes sense:
\[
\text{MSE}=\frac{1}{N}\sum_{i=1}^{N}(y_i-\hat{y}_i)^2.
\]
Classification predicts probabilities; the standard loss is cross entropy.

\subsection{Binary cross entropy (BCE)}
For $y_i\in\{0,1\}$ and predicted probability $p_i$:
\[
\mathcal{L}_{\text{BCE}}(y,p)=
-\frac{1}{N}\sum_{i=1}^{N}\left[y_i\log(p_i)+(1-y_i)\log(1-p_i)\right].
\]

\section{The PyTorch Training Loop: What \texttt{fit} Is Hiding}

Gradient-based training repeats the same core steps each epoch:

\begin{tcolorbox}[colframe=orange!70!black,colback=orange!4,title=\textbf{\color{orange!80!black}One epoch = one parameter update cycle}]
\begin{enumerate}[leftmargin=*]
  \item Zero old gradients
  \item Forward pass: compute outputs
  \item Compute loss
  \item Backward pass: compute gradients
  \item Update parameters
\end{enumerate}
\end{tcolorbox}

Mathematically, gradient descent updates:
\[
w \leftarrow w - \eta\nabla_w\mathcal{L},
\qquad
b \leftarrow b - \eta\nabla_b\mathcal{L},
\]
where $\eta$ is the learning rate.

\section{Tutorial: Breast Cancer Classification (End-to-End)}

\subsection{Step 1: Load, split, and standardize}
\begin{tcolorbox}[colframe=black!40,colback=black!2,title=\textbf{Why standardization matters}]
If one feature ranges near $10^6$ and another near $10^{-3}$, gradient-based training becomes unstable and slow.
Standardization makes features comparable.
\end{tcolorbox}

\begin{lstlisting}[language=Python,caption={Load data, train/test split, and standardize (scikit-learn)}]
import numpy as np
import torch
import torch.nn as nn

from sklearn.datasets import load_breast_cancer
from sklearn.model_selection import train_test_split
from sklearn.preprocessing import StandardScaler

# Load dataset
data = load_breast_cancer()
X = data.data          # shape: (N, D)
y = data.target        # shape: (N,)

# Split
X_train, X_test, y_train, y_test = train_test_split(
    X, y, test_size=0.33, random_state=123
)

# Standardize using train statistics only
scaler = StandardScaler()
X_train = scaler.fit_transform(X_train)
X_test  = scaler.transform(X_test)
\end{lstlisting}

\subsection{Step 2: Convert to PyTorch tensors (watch shapes!)}
PyTorch often expects targets shaped as $(N,1)$ for binary losses.

\begin{lstlisting}[language=Python,caption={Convert numpy arrays to torch tensors (float32) and reshape targets}]
X_train_t = torch.from_numpy(X_train).float()
X_test_t  = torch.from_numpy(X_test).float()

y_train_t = torch.from_numpy(y_train).float().view(-1, 1)  # (N,1)
y_test_t  = torch.from_numpy(y_test).float().view(-1, 1)   # (N,1)

D = X_train.shape[1]
\end{lstlisting}

\section{Two Correct Implementations (Choose One)}

\subsection{Option A: Linear + Sigmoid + BCELoss (simple, but less stable)}

\begin{tcolorbox}[colframe=red!60!black,colback=red!3,title=\textbf{\color{red!70!black}Note}]
This is correct, but can be numerically unstable due to exponentials and logs.
Option B is usually preferred in practice.
\end{tcolorbox}

\begin{lstlisting}[language=Python,caption={Model A: sigmoid inside the model + BCELoss}]
model = nn.Sequential(
    nn.Linear(D, 1),
    nn.Sigmoid()
)

criterion = nn.BCELoss()
optimizer = torch.optim.Adam(model.parameters(), lr=1e-3)
\end{lstlisting}

\subsubsection{Training loop (same structure always)}
\begin{lstlisting}[language=Python,caption={Training loop with train/test loss tracking (Model A)}]
n_epochs = 200
train_losses, test_losses = [], []

for epoch in range(n_epochs):
    # ---- TRAIN ----
    optimizer.zero_grad()
    p_train = model(X_train_t)                    # probabilities (0,1)
    loss_train = criterion(p_train, y_train_t)
    loss_train.backward()
    optimizer.step()

    # ---- TEST LOSS ----
    with torch.no_grad():
        p_test = model(X_test_t)
        loss_test = criterion(p_test, y_test_t)

    train_losses.append(loss_train.item())
    test_losses.append(loss_test.item())

print("Final train loss:", train_losses[-1])
print("Final test  loss:", test_losses[-1])
\end{lstlisting}

\subsubsection{Accuracy}
\begin{lstlisting}[language=Python,caption={Accuracy from probabilities}]
def accuracy_from_prob(p, y_true):
    # p: (N,1) probabilities
    y_pred = (p >= 0.5).float()
    return (y_pred.eq(y_true)).float().mean().item()

with torch.no_grad():
    p_train = model(X_train_t)
    p_test  = model(X_test_t)

print("Train accuracy:", accuracy_from_prob(p_train, y_train_t))
print("Test accuracy: ", accuracy_from_prob(p_test,  y_test_t))
\end{lstlisting}

\subsection{Option B (Recommended): Linear logits + BCEWithLogitsLoss (stable)}

\begin{tcolorbox}[colframe=green!60!black,colback=green!3,title=\textbf{\color{green!50!black}Recommended in practice}]
Use logits directly and let the loss handle the sigmoid internally in a numerically stable way.
\end{tcolorbox}

\begin{lstlisting}[language=Python,caption={Model B: logits output + BCEWithLogitsLoss}]
model = nn.Linear(D, 1)                # outputs logits (any real number)
criterion = nn.BCEWithLogitsLoss()     # stable sigmoid + BCE
optimizer = torch.optim.Adam(model.parameters(), lr=1e-3)
\end{lstlisting}

\subsubsection{Training loop (identical skeleton)}
\begin{lstlisting}[language=Python,caption={Training loop with logits (Model B)}]
n_epochs = 200
train_losses, test_losses = [], []

for epoch in range(n_epochs):
    optimizer.zero_grad()
    logits_train = model(X_train_t)                 # logits
    loss_train = criterion(logits_train, y_train_t)
    loss_train.backward()
    optimizer.step()

    with torch.no_grad():
        logits_test = model(X_test_t)
        loss_test = criterion(logits_test, y_test_t)

    train_losses.append(loss_train.item())
    test_losses.append(loss_test.item())

print("Final train loss:", train_losses[-1])
print("Final test  loss:", test_losses[-1])
\end{lstlisting}

\subsubsection{Accuracy from logits (no sigmoid needed)}
Because $\sigma(a)\ge 0.5 \Leftrightarrow a\ge 0$, we classify by sign:
\[
\hat{y} = \mathbb{1}[a\ge 0].
\]

\begin{lstlisting}[language=Python,caption={Accuracy from logits}]
def accuracy_from_logits(logits, y_true):
    y_pred = (logits >= 0).float()     # sign rule
    return (y_pred.eq(y_true)).float().mean().item()

with torch.no_grad():
    logits_train = model(X_train_t)
    logits_test  = model(X_test_t)

print("Train accuracy:", accuracy_from_logits(logits_train, y_train_t))
print("Test accuracy: ", accuracy_from_logits(logits_test,  y_test_t))
\end{lstlisting}

\section{Saving and Loading the Model}

\subsection{What is a \texttt{state\_dict}?}
\texttt{state\_dict} is an ordered dictionary of parameters (weights, biases, etc.).
For a linear layer, it contains:
\[
W \in \mathbb{R}^{1\times D},\qquad b \in \mathbb{R}.
\]

\subsection{Save}
\begin{lstlisting}[language=Python,caption={Save learned parameters}]
torch.save(model.state_dict(), "mymodel.pt")
\end{lstlisting}

\subsection{Load (recreate architecture first)}
\begin{lstlisting}[language=Python,caption={Load parameters into a new model with same architecture}]
model2 = nn.Linear(D, 1)
model2.load_state_dict(torch.load("mymodel.pt"))
model2.eval()
\end{lstlisting}

\subsection{Verify loaded model}
\begin{lstlisting}[language=Python,caption={Check test accuracy of loaded model}]
with torch.no_grad():
    logits_test_2 = model2(X_test_t)

print("Reloaded model test accuracy:",
      accuracy_from_logits(logits_test_2, y_test_t))
\end{lstlisting}

\section{Train vs Validation vs Test (Student-Friendly Explanation)}

\begin{tcolorbox}[colframe=blue!60!black,colback=blue!3,title=\textbf{\color{blue!70!black}Why splitting matters}]
A model can look great on training data but fail on new data. \\
We split data to estimate real-world performance and to avoid overfitting.
\end{tcolorbox}

\subsection{Two-way split}
\begin{itemize}[leftmargin=*]
  \item \textbf{Train set}: used to learn parameters $(w,b)$.
  \item \textbf{Test set}: used to estimate performance on unseen data.
\end{itemize}

\subsection{Three-way split}
\begin{itemize}[leftmargin=*]
  \item \textbf{Train}: learn parameters.
  \item \textbf{Validation}: tune hyperparameters (learning rate, epochs, architecture).
  \item \textbf{Test}: final unbiased evaluation (use once at the end).
\end{itemize}

\subsection{Parameters vs hyperparameters}
\begin{tcolorbox}[colframe=purple!60!black,colback=purple!3,title=\textbf{\color{purple!70!black}Important distinction}]
\begin{itemize}[leftmargin=*]
  \item \textbf{Parameters} (learned): $w, b$
  \item \textbf{Hyperparameters} (chosen by you): learning rate $\eta$, optimizer, number of epochs, model size
\end{itemize}
Validation helps you choose hyperparameters.
\end{tcolorbox}

\section{Exercises (to make it stick)}

\subsection*{\color{blue!70!black}Exercise 1: Track accuracy per epoch}
During training, compute train and test accuracy each epoch and store them in arrays. Plot:
\[
\text{loss vs epoch}, \qquad \text{accuracy vs epoch}.
\]

\subsection*{\color{blue!70!black}Exercise 2: Confirm two prediction rules match}
Verify these produce the same predictions:
\[
\hat{y}=\mathbb{1}[a\ge 0]
\quad \text{and} \quad
\hat{y}=\mathbb{1}[\sigma(a)\ge 0.5].
\]

\subsection*{\color{blue!70!black}Exercise 3: Learning rate experimentation}
Try $\eta \in \{10^{-1}, 10^{-2}, 10^{-3}, 10^{-4}\}$ and observe:
\begin{itemize}[leftmargin=*]
  \item too high $\Rightarrow$ loss explodes / diverges
  \item too low $\Rightarrow$ training is very slow
  \item just right $\Rightarrow$ smooth decreasing loss
\end{itemize}

\section{Chapter Summary (One Page)}
\begin{tcolorbox}[colframe=black!50,colback=black!3,title=\textbf{\color{black}You should now understand...}]
\begin{itemize}[leftmargin=*]
  \item Linear classification uses $a=w^\top x+b$ and separates classes by the sign of $a$.
  \item Logistic regression converts $a$ to probability with $\sigma(a)=\frac{1}{1+e^{-a}}$.
  \item Classification uses BCE-style losses; \texttt{BCEWithLogitsLoss} is numerically stable.
  \item The PyTorch training loop is: zero grad $\to$ forward $\to$ loss $\to$ backward $\to$ step.
  \item Accuracy is computed by comparing predicted labels with true labels.
  \item Models are saved with \texttt{state\_dict} and loaded by recreating architecture + loading weights.
  \item Train/validation/test splits support generalization and hyperparameter tuning.
\end{itemize}
\end{tcolorbox}